\medskip
\noindent\emph{Comparison with the variable-shape route.}
The other construction named in the supplied GALOIS F.5--F.6 and
PATHS D.14--D.16 varies the indexing poset and retains the groupoids
$\Pi_1\operatorname{Map}(BJ,X)$. The stated class of tests consists
of posets of finite chain dimension, not merely finite posets;
in particular it contains infinite discrete sets. This is distinct
from increasing the coefficient truncation in E.1--E.4.
Even on finite tests it cannot be constructed from $\Pi_1X$ alone.
Take $J$ to be the face poset of the boundary of a tetrahedron, so
$BJ\simeq S^2$, and take $X=K(\mathbb Z,2)$. Then
\[
 \pi_0\operatorname{Map}(BJ,X)
       =[S^2,K(\mathbb Z,2)]
       =H^2(S^2;\mathbb Z)=\mathbb Z,
 \qquad \Pi_1X\simeq *.
\]
Thus $\Pi_1\operatorname{Map}(BJ,X)$ has infinitely many components,
whereas $\operatorname{Map}(BJ,\Pi_1X)$ is contractible.
This proves the needed distinction directly. No reconstruction
theorem from the preliminary variable-shape chapters is used as
an input to the proofs of this return.
